% Lösungen Einheit 2 von 5 – Quader und Würfel (zusammenbau v0.1)
\einheitenkopf{Einheit 2 von 5 · Quader und Würfel}

%% TODO Lösung der Erklärzeile 14g) fehlt in der Bank
\erg{14}{a) Volumen \quad b) $6 \cdot 3 \cdot 2 = 36$ cm³ \quad c) $7 \cdot 5 \cdot 2 = 70$ cm³ \quad d) $8 \cdot 5 \cdot 3 = 120$ m³ \quad e) $9 \cdot 4 \cdot 3 = 108$ dm³ \quad f) $5 \cdot 5 \cdot 3 = 75$ cm³ \quad g) \ldots \quad h) $5 \cdot 5 \cdot 5 = 125$ cm³ \quad i) $4\,500 : 1\,000 = 4{,}5$ l \quad j) $2 \cdot (7 \cdot 3 + 7 \cdot 2 + 3 \cdot 2) = 82$ cm² \quad k) $6 \cdot 5 \cdot 5 = 150$ dm² \quad l) $90 : (5 \cdot 3) = 6$ cm \quad m) $6 \cdot 4 \cdot 2 + 3 \cdot 4 \cdot 2 = 48 + 24 = 72$ dm³ \quad n) $8 \cdot 8 \cdot 8 = 512$ cm³ (nicht $8 \cdot 8 = 64$, nicht die Oberfläche $384$ cm²)}
\erg{15}{a) $8 \cdot 5 + 2 \cdot 8 \cdot 4 + 2 \cdot 5 \cdot 4 = 144$ dm²}
\erg{16}{a) $30$ cm, denn $30 \cdot 30 \cdot 30 = 27\,000$}
\erg{17}{a) $3 \cdot 1\,000 = 3\,000$ l}
\erg{18}{a) $50$ l $= 50\,000$ cm³; $50\,000 : (50 \cdot 25) = 40$ cm}
\erg{19}{a) Nils hat das Volumen gerechnet. Oberfläche: $2 \cdot (14 + 35 + 10) = 118$ cm² \quad b) Länge mal Breite ist die Zahl der Würfel in der untersten Schicht; es gibt so viele Schichten, wie der Quader in cm hoch ist. Also Länge mal Breite mal Höhe Würfel. \quad c) $62 \cdot 34 \cdot 36 = 75\,888$ cm³ $\approx 75{,}9$ l – ja, etwa $75$ l \quad d) je zweimal $14$ cm², $7$ cm² und $2$ cm²; zusammen $46$ cm²}

19d)

\netzquader{3.5}{1}{0.5}{$7$ cm}{$2$ cm}{$1$ cm}

